\chapter{Formulario de referencia rápida}
\label{anexo:formulario}

Este anexo recopila, en un solo lugar, todas las ecuaciones marcadas como \textcolor{ucred}{\textbf{Ecuación importante}} a lo largo del libro, en el orden en que aparecen. Sirve como formulario de repaso rápido antes de una evaluación; para el contexto completo y la derivación de cada resultado, siga el hipervínculo hasta la sección correspondiente.

\tcblistof[\chapter*]{ecnimportante}{Índice de ecuaciones importantes}

\bigskip
\noindent\textbf{Referencias recomendadas.} Los apuntes de este libro siguen la formulación clásica del método de los elementos finitos presentada, entre otros, en:
\begin{itemize}[leftmargin=1.4em]
  \item K.-J.\ Bathe, \emph{Finite Element Procedures}, Prentice Hall.
  \item O.\,C.\ Zienkiewicz, R.\,L.\ Taylor, \emph{The Finite Element Method}, Butterworth-Heinemann.
  \item R.\,D.\ Cook, D.\,S.\ Malkus, M.\,E.\ Plesha, R.\,J.\ Witt, \emph{Concepts and Applications of Finite Element Analysis}, Wiley.
  \item T.\,J.\,R.\ Hughes, \emph{The Finite Element Method: Linear Static and Dynamic Finite Element Analysis}, Dover.
  \item B.\ Irons, S.\ Ahmad, \emph{Techniques of Finite Elements}, Ellis Horwood (origen del \emph{patch test}).
\end{itemize}
